\section{Chest X-Ray}


\section*{Patient focus}
This chapter explains the examination or clinical problem in accessible language. Patients should use it to prepare questions and understand the purpose, experience, benefits, and limitations of imaging; local instructions and the official report take priority.


\subsubsection*{What Is This Test?}
A chest X-ray uses a small amount of radiation to create images of your heart, lungs, airways, blood vessels, and chest bones. It's quick, inexpensive, and often the first imaging test for chest symptoms.

\subsubsection*{Alternative Names}
CXR, chest radiograph, chest X-ray, lung X-ray, thoracic X-ray


\subsubsection*{Principle}
Chest X-ray uses the same radiographic principle as abdominal X-ray. The posteroanterior (PA) and lateral views are standard. The X-ray beam is attenuated by the chest structures: the heart, mediastinum, and diaphragm appear radiopaque (white); the lungs, being air-filled, appear radiolucent (black).

\subsubsection*{History}
The chest X-ray was among the first medical imaging techniques, used within months of Röntgen's 1895 discovery of X-rays. Mass chest X-ray screening for tuberculosis was implemented in the 1930s and continued for decades.

\subsubsection*{Why This Test Is Ordered}
Your doctor orders a chest X-ray when you have a cough, fever, shortness of breath, chest pain, or after a trauma to your chest. It answers the question: is there pneumonia, fluid around your lungs, a collapsed lung, an enlarged heart, or something else wrong in your chest?

\subsubsection*{Is This a Primary or Secondary Test}
A primary test. One of the most commonly ordered imaging tests and usually the first look at the chest when symptoms arise.

\subsubsection*{What a chest radiograph can miss}
Subtle or early disease may be missed, especially on a single portable AP image. Sensitivity should be linked to the suspected finding and view conditions; ongoing clinical concern may warrant further evaluation.

\subsubsection*{How This Test Is Performed}
You stand in front of an X-ray machine with your chest against a detector. You take a deep breath and hold it while the X-ray is taken. Two images are usually taken: one from the front and one from the side. The test takes about five minutes. There is no pain involved.

\subsubsection*{What Preparation Is Needed}
Remove any jewelry, glasses, or metal objects from your chest area. Tell your technician if you are pregnant, although the radiation dose from a chest X-ray is very low. Women should inform staff of any possibility of pregnancy.

\subsubsection*{Contraindications and Precautions}
Pregnancy is not an automatic contraindication; the examination should be justified and optimized, and alternatives considered when appropriate. For contrast-enhanced studies, review the indication, prior severe reactions, and renal function according to local protocol.

\subsubsection*{Risks}
Risks associated with X-ray imaging include:
\begin{itemize}
    \item \textbf{Radiation exposure}: Exposure to low levels of ionizing radiation, with cumulative lifetime dose optimization guided by the ALARA (As Low As Reasonably Achievable) principle.
\end{itemize}


A chest X-ray is completely non-invasive. You can resume all normal activities immediately. Results are typically available within thirty minutes to two hours.

\subsubsection*{Understanding Your Results}

\subsubsection*{Normal Results}
A normal result shows clear lungs with no infiltrates, masses, fluid, or pneumothorax. Normal heart size and mediastinal contour.

\subsubsection*{Follow-up Tests}
If the chest X-ray is abnormal, your doctor may order a CT scan of the chest for more detail, blood tests to evaluate infection or inflammation, or a follow-up X-ray after treatment to confirm resolution.
\clearpage
